% =============================================================================
% OBSERVATION LOG
% Status: ACTIVE - Updated throughout research process
% =============================================================================

% This appendix maintains a chronological log of observations made during
% the research process. Observations are tagged with dates and linked to
% relevant figures/sections where applicable.

\subsection*{Log Format}
Each observation entry includes:
\begin{itemize}
    \item \textbf{Date}: When the observation was made
    \item \textbf{Category}: Type of observation (data, method, result, idea)
    \item \textbf{Description}: The observation itself
    \item \textbf{Link}: Reference to relevant section, figure, or data
\end{itemize}

\hrulefill

\subsection*{2026-01-13}

\begin{description}
    \item[OBS-001] \textbf{[Data]} Dataset loaded successfully: 97 DJI drone images from Santa Ines field, ~10-12 MB each.

    \item[OBS-002] \textbf{[Result]} UMAP projection reveals clear bimodal structure---main cluster of 89 images and outlier cluster of 8 images at bottom-right of embedding space.

    \item[OBS-003] \textbf{[Result]} Outlier images (DJI\_0469, DJI\_0509, DJI\_0552-0557) visually confirmed as high-altitude overview shots, distinct from close-up crop inspection images.

    \item[OBS-004] \textbf{[Method]} CLIP embeddings capture altitude/scale differences effectively---semantic clustering emerges without altitude metadata being provided.
\end{description}

\subsection*{2026-01-14}

\begin{description}
    \item[OBS-005] \textbf{[Result]} Similarity index computed successfully. Query tests confirm semantically consistent retrieval---high-altitude queries return only high-altitude neighbors.

    \item[OBS-006] \textbf{[Result]} Network analysis: 97.3\% of k-NN edges stay within category. Only 8/291 edges cross category boundary.

    \item[OBS-007] \textbf{[Result]} Cluster cohesion asymmetry: high-altitude cluster is 35\% tighter (mean neighbor distance 0.276 vs 0.424). May indicate greater visual homogeneity of overview shots.

    \item[OBS-008] \textbf{[Idea]} Cohesion metric could serve as diversity indicator---tight clusters = homogeneous content, diffuse clusters = varied content. Useful for coverage assessment.

    \item[OBS-009] \textbf{[Method]} R/ggplot2 workflow established for reproducible analysis. Scripts: \texttt{analysis.R}, \texttt{similarity\_analysis.R}.
\end{description}

\hrulefill

\subsection*{Template for Future Entries}

\begin{description}
    \item[OBS-XXX] \textbf{[Category]} Description of observation. \textit{See Section~X / Figure~X}
\end{description}

% =============================================================================
% OBSERVATION INDEX BY CATEGORY
% =============================================================================

\subsection*{Index by Category}

\textbf{Data observations:} OBS-001

\textbf{Method observations:} OBS-004, OBS-009

\textbf{Result observations:} OBS-002, OBS-003, OBS-005, OBS-006, OBS-007

\textbf{Ideas/hypotheses:} OBS-008
